\chapter{流体动力学基本方程}

\begin{yijubox}
E1 \texttt{8-1考情分析.pptx} slide33：第四章流体动力学基础＝\textbf{考试计算题重点}。\\
E2 \texttt{9-1考点分析.pptx}：计算题占 50\% 分值；\textbf{第四章是主要考点，需掌握计算套路}。\\
E1 slide16：计算题考\textbf{伯努利方程、动量方程}等。\\
\textbf{本章结论}：这是全书分值最高的计算章节之一。按试卷结构（计算题 3 题共 58 分，
其中必然有 1--2 题落在这里，约 15--20 分），再加选择、填空、判断各 1--2 题，
本章合计约 \textbf{20--30 分}。复习策略＝\textbf{把 4.2 伯努利方程与 4.3 动量方程的两套解题套路练到条件反射}，
4.1 只背公式与适用条件，4.4 只记公式形状。
\end{yijubox}

\section{理想流体的运动微分方程}

\begin{kaodian}{3}{理想流体的特征与欧拉运动微分方程}
\textbf{是什么：}理想流体＝\textbf{无黏性}的假想流体。动力学分析总是"先理想、后实际"：
先在 $\mu=0$ 的条件下建立基本方程，再把黏性作为修正项加回去（加上黏性项就得到第 8 章的纳维--斯托克斯方程）。

\begin{gongshi}{欧拉运动微分方程（分量式，式 4.1）}
\[
\begin{aligned}
X-\frac{1}{\rho}\frac{\partial p}{\partial x}&=\frac{\dd u_x}{\dd t}\\
Y-\frac{1}{\rho}\frac{\partial p}{\partial y}&=\frac{\dd u_y}{\dd t}\\
Z-\frac{1}{\rho}\frac{\partial p}{\partial z}&=\frac{\dd u_z}{\dd t}
\end{aligned}
\]
矢量形式（$f$ 为单位质量力，$\bu=(u_x,u_y,u_z)$）：
\[ \bm{f}-\frac{1}{\rho}\grad p=\frac{\dd \bu}{\dd t} \]
\textbf{含义：}对理想流体中的流体微团直接应用牛顿第二定律。左端第一项＝单位质量力（重力、惯性力等），
左端第二项＝\textbf{单位质量表面力的合力}（理想流体表面力只有法向压力，没有切应力，所以只剩 $-\grad p/\rho$），
右端＝微团的加速度。三项关系也可写成\textbf{单位质量的质量力＋单位质量的压力＋单位质量的惯性力＝0}。\\
\textbf{适用条件：}\textbf{理想流体}（唯一限制）；不要求定常、不要求不可压缩、不要求无旋。
所以它是"适用于所有理想流体流动"的普遍方程。实际流体必须补上黏性项。\\
\textbf{常考结论：}"理想流体的表面力只有压力没有切应力"——因为 $\mu=0$，由牛顿内摩擦定律 $\tau=\mu\,\dd u/\dd y$ 直接得 $\tau=0$。
\end{gongshi}

\begin{gongshi}{兰姆--葛罗米柯（Lamb--Gr\"om\'eko）形式，式 4.4}
把加速度按矢量恒等式展开：
\[ \frac{\dd \bu}{\dd t}=\frac{\partial \bu}{\partial t}+\grad\!\left(\frac{u^2}{2}\right)-\bu\times\bomega,\qquad \bomega=\nabla\times\bu \]
再设质量力有势（$f=-\grad W$），欧拉方程化为
\[ \grad\!\left(W+\frac{p}{\rho}+\frac{u^2}{2}\right)=\bu\times\bomega-\frac{\partial \bu}{\partial t} \]
\textbf{含义：}把加速度拆成"局部加速度＋位变加速度"两部分，后者又分成一个\textbf{标量梯度项} $\grad(u^2/2)$
和一个\textbf{旋涡项} $\bu\times\bomega$。好处是：右端的旋涡项与压强无关，
一旦流动\textbf{无旋}（$\bomega=0$）右端就只剩局部项，方程变成纯梯度形式，可直接积分得到伯努利积分——
这正是下面两个积分式的来源。\\
\textbf{适用条件：}理想流体；质量力有势。定常流时 $\partial\bu/\partial t=0$，右端只剩 $\bu\times\bomega$。\\
\textbf{考试口径：}该形式只需认得出、说得出"旋涡项"，\textbf{不要求推导}。
\end{gongshi}

\textbf{常考题型：}选择/判断（理想流体的表面力有哪两种（答：只有压力）；
"欧拉运动微分方程适用于哪些流动"（答：所有理想流体流动，不限该流动是否定常、可压缩、无旋））；
简答（写出欧拉方程的分量式并说明各项物理意义）。
\textbf{重要度 \xstarfull{3}}\quad\yiju{E1 slide33 第四章＝考试计算题重点，但针对 4.1 只要求基础概念与公式；E2 第四章"需掌握计算套路"}\quad\nandu{易}

\begin{banfa}
\textbf{①"三项平衡"记法}：欧拉方程就是\textbf{单位质量}的牛顿第二定律——
\[ \underbrace{f}_{\text{质量力}}-\underbrace{\frac{1}{\rho}\grad p}_{\text{压力}}- \underbrace{\frac{\dd\bu}{\dd t}}_{\text{惯性力}}=0 \]
三种力一张口就能背，不必硬记偏导数形式。\\
\textbf{②分量式速记}：$x$ 方向就是 $X-\dfrac{1}{\rho}\dfrac{\partial p}{\partial x}=\dfrac{\dd u_x}{\dd t}$，
$y$、$z$ 只是把字母换掉。注意左端对 $p$ 是\textbf{偏导数}（且带负号），
右端速度对时间是\textbf{全导数/随体导数}（第 3 章 3.1.3 的结论）——这一处是填空题最爱挖的坑。\\
\textbf{③它和"伯努利积分"的关系}：欧拉方程是\textbf{微分}形式、普遍适用；
伯努利积分是它在"不可压缩＋定常＋质量力只有重力＋沿流线"下的\textbf{一次积分}，条件多得多。
选择题问"下列哪个方程适用范围最广"，答\textbf{欧拉运动微分方程}。
\end{banfa}
\end{kaodian}

\begin{kaodian}{3}{理想流体定常流的伯努利积分（含适用条件）}
\textbf{是什么：}把欧拉方程在特定条件下沿流线积分一次，得到的一条机械能守恒关系式。
定常、不可压缩、质量力只有重力时，积分常数为 $C$。

\begin{gongshi}{理想流体元流（微小流束）的伯努利方程，式 4.5、4.6}
沿流线积分：
\[ z+\frac{p}{\rho g}+\frac{u^2}{2g}=C \]
同一条流线上取任意两点 1、2：
\[ z_1+\frac{p_1}{\rho g}+\frac{u_1^2}{2g}=z_2+\frac{p_2}{\rho g}+\frac{u_2^2}{2g} \]
\textbf{含义：}单位重力流体沿流线的三种机械能（位能、压能、动能，分别以长度量纲的水头表示）\textbf{总量守恒}；
$C$ 称\textbf{伯努利常数}，同一流线上 $C$ 相同。\\
\textbf{适用条件（\textbf{四条，缺一不可}，必背）：}
\begin{xiaowen}
  \item \textbf{理想流体}（无黏，$\mu=0$），因此方程就是能量守恒，不含水头损失项；
  \item \textbf{不可压缩}流体，$\rho$ 为常数；
  \item \textbf{定常流动}（稳定流动），质点沿流线方向运动；
  \item 质量力\textbf{只有重力}一项。\\
\end{xiaowen}
\textbf{最关键的范围限制，见下表。}
\end{gongshi}

\begin{center}
\begin{tabularx}{\textwidth}{@{}lXX@{}}
\toprule
积分名称 & 成立范围 & 能否跨流线用 \\
\midrule
\textbf{伯努利积分} & 不可压缩、定常、质量力只有重力、\textbf{沿同一条流线}（此时流动可以有旋） & \textbf{不能}。有旋时 $C$ 逐流线不同 \\
\textbf{拉格朗日积分}（无旋流动积分） & 在伯努利积分的条件上再加\textbf{无旋}（$\bomega=0$，全流场为势流） & \textbf{能}。此时全流场 $C$ 相同，可跨流线取点 \\
\bottomrule
\end{tabularx}
\end{center}

\textbf{几何意义与物理意义（简答必答）：}
\textbf{几何}＝位置水头 $z$ ＋ 压强水头 $p/(\rho g)$ ＋ 流速水头 $u^2/(2g)$ ＝ 总水头（三者都是长度量纲，单位 m）；
\textbf{物理}＝单位重量流体的位能 ＋ 压能 ＋ 动能 ＝ 总机械能（单位 m，即 $\mathrm{J/N}$）。
换言之：理想不可压缩流体在重力场中定常流动时，\textbf{单位重量流体的总机械能沿流线守恒}。

\textbf{常考题型：}填空（说出理想流体伯努利方程中三项的物理意义，或由上、下游读出缺少的一项）；
判断（"理想流体伯努利方程对任意流动都成立"$\to$ 错，必须满足四条）；
简答（写出适用条件并说明无黏流为何不含水头损失项）。
\textbf{重要度 \xstarfull{3}}\quad\yiju{E1 slide16 填空考"公式含义"；E2 第四章"需掌握计算套路"；E6 教材 4.1 节为该积分式专节}\quad\nandu{易}

\begin{banfa}
\textbf{①"沿流线还是全场"是本章最贵的 3 分}：看到题目里出现"有旋/旋涡/剪切"，
理想流体伯努利方程\textbf{只能在同一条流线上取两点}；只有题目明确是"\textbf{势流/无旋}"才能跨流线。
考场速判句：\textbf{"无黏管沿流线，无旋才全场"}。\\
\textbf{②四条适用条件的"漏项自检"}：定常？不可压缩？无黏（理想）？只有重力？
把四个问号在心里过一遍，再看题干有没有给水泵/水轮机（给了就是 4.2 的带泵方程，不能只用 4.5 式）。\\
\textbf{③伯努利积分与拉格朗日积分不要搞混}：记法是"\textbf{拉格朗日管全场（无旋），伯努利只管一条线}",
对应的物理图像分别是"势流"与"流线"。
\end{banfa}
\end{kaodian}

\section{伯努利方程}

\begin{kaodian}{5}{实际流体总流的伯努利方程（含动能修正系数与水头损失）}
\textbf{是什么：}把元流（流线）的伯努利方程沿总流断面"平均"之后得到的工程实用式。
两条修正：① 用断面\textbf{平均流速} $v$ 代替点流速 $u$，需乘\textbf{动能修正系数} $\alpha$；
② 实际流体有黏性，沿程必然损失机械能，需减去\textbf{水头损失} $h_w$。

\begin{gongshi}{实际流体总流伯努利方程}
\[ z_1+\frac{p_1}{\rho g}+\frac{\alpha_1 v_1^2}{2g}=z_2+\frac{p_2}{\rho g}+\frac{\alpha_2 v_2^2}{2g}+h_w \]
断面上三种水头之和 $z+\dfrac{p}{\rho g}+\dfrac{\alpha v^2}{2g}$ 称该断面的\textbf{总水头} $H_p$，
两断面总水头之差就是两断面之间的水头损失；水头损失由沿程损失与局部损失两部分组成：
\[ H_{p1}-H_{p2}=h_w,\qquad h_w=h_f+h_j,\qquad i=\frac{h_w}{L} \]
（$h_f$＝沿程水头损失，$h_j$＝局部水头损失，$i$＝水力坡度＝单位管长的水头损失，即\textbf{总水头线的斜率}。）
\textbf{含义：}上游断面的总水头＝下游断面的总水头＋两断面之间\textbf{单位重量流体}损失掉的机械能。
$z$＝位置水头（位能），$p/(\rho g)$＝压强水头（压能），$\alpha v^2/(2g)$＝按平均流速算的流速水头（动能），
$h_w$＝水头损失（单位 m 水柱），本质是\textbf{黏性耗散}，机械能变成了热能，\emph{不可逆}。\\
\textbf{适用条件（五条，必背）：}
\begin{xiaowen}
  \item \textbf{不可压缩}流体（$\rho=$ 常数；气体低速流动也可用，但要取平均密度）；
  \item \textbf{定常流动}；
  \item 作用在流体上的质量力\textbf{只有重力}；
  \item 所取的两个断面都必须是\textbf{渐变流（缓变流）断面}——因为只有渐变流断面上压强才按\textbf{静压分布}
        $z+p/(\rho g)=$ 常数，否则 $p$ 与断面平均流速没有对应的能量含义。
        渐变流＝流线之间夹角小、曲率半径大的流动；急变流＝流线不平行的直线且夹角大，或曲率半径很小。
        注意：\textbf{两断面之间的流动可以是急变流}（含局部水头损失），只是\emph{断面本身}要取在渐变流处；
  \item 两断面间\textbf{沿程无流量分出或汇入}（有分流/汇流要分段列方程，或按 4.3 的多进出口动量方程处理），
        且两断面之间除水头损失外\textbf{无其他能量输入输出}（有则用下一考点的带泵/带水轮机式）。
\end{xiaowen}
\textbf{位置水头的基准}：$z$ 取同一水平基准面为准。基准面可任取，但\textbf{两个断面必须用同一个基准面}；
工程上习惯取较低断面的中心为基准面，使 $z_1=0$、$z_2>0$，少一笔负数运算。\\
\textbf{动能修正系数 $\alpha$：}同一断面上按点流速算的实际动能与按平均流速算的假想动能之比，
$\alpha=\dfrac{\displaystyle\int_A u^3\,\dd A}{v^3 A}\ \ge 1$，反映\textbf{断面速度分布的不均匀程度}——越不均匀 $\alpha$ 越大。
层流（抛物线型速度分布，分布最不均匀）$\alpha=2$；湍流 $\alpha\approx 1$（严格为 $1.05\sim1.1$）；
工程计算中因流速水头占总水头比例小，\textbf{题目不特别说明就取 $\alpha=1$}。\\
\textbf{水头损失的写法：}教材用 $h'_{l-2}$ 表示沿微小流束两断面间的损失，用 $h_{l-2}$（本讲义统一写作 $h_w$）表示总流两断面间的损失，
二者区别只在"元流/总流"，数值含义相同（$\mathrm{m}$ 水柱）。
\end{gongshi}

\begin{tishi}
\textbf{理想流体是实际流体的特例}：令 $h_w=0$、$\alpha=1$，实际流体总流伯努利方程就退化成理想流体元流伯努利方程。
反过来，\textbf{判断题里"实际流体伯努利方程就是理想流体伯努利方程加上水头损失"是对的}，
但"\textbf{两式可以直接互相替代}"是\textbf{错}的——断面上 $u$ 与 $v$、$\alpha$ 的差别不能忽略。
\end{tishi}

\textbf{常考题型：}填空（"水头损失的物理意义是（单位重力流体损失的机械能）"）；
选择（2022 年 818 单选第 6 题考"应用伯努利方程所选的计算点位于何处"）；
简答（写出适用条件；说明 $\alpha$ 的物理意义与取值）；计算（几乎本章所有计算题的第一步）。
\textbf{重要度 \xstarfull{5}}\quad\yiju{E4 2022 年 818 单选第 6 题（伯努利方程所选计算点位置）；E1 slide33 第四章＝考试计算题重点；E1 slide16 计算考伯努利方程}\quad\nandu{中}

\begin{zhenti}{2022 年 818 第 6 题}{单项选择题（回忆版）}
应用伯努利方程时，所选计算点应位于何处？\\
(A) 任意位置\quad (B) 断面中心点\quad (C) 渐变流断面上（且使方程未知量最少处）\quad (D) 管壁处
\end{zhenti}
\noindent\textbf{答案：(C)。}两个断面必须是\textbf{渐变流断面}；断面上的计算点习惯取\textbf{断面中心}，
因为那里最容易用测压管高度直接读出 $z+p/(\rho g)$。选点的唯一目标是\textbf{让方程里的未知数最少}。

\begin{yicuo}
\textbf{①基准面不统一}：位置水头 $z$ 的基准面两个断面必须相同，且\textbf{压强基准也必须两个断面一致}——
常用相对压强（表压）配相对压强。\textbf{一个断面用绝对压强、另一个用相对压强}是计算题最经典的失分点。\\
\textbf{②把计算点取在急变流断面}（弯头内部、闸门后、孔口出口紧上游）：那里断面上压强\textbf{不按静压分布}，
$z+p/(\rho g)=$ 常数的前提已破坏，方程不成立。\\
\textbf{③把 $h_w$ 忘掉或加错端}：实际流体中能量只会减少，$h_w$ 一定写在\textbf{下游一侧}（本式为 $+h_w$）。
判据：理想流体两断面总水头相等；实际流体下游总水头更低。\\
\textbf{④有泵/水轮机却仍用无泵式}：见下一考点。
\end{yicuo}
\end{kaodian}

\begin{kaodian}{5}{带泵（或水轮机）的伯努利方程与泵的功率}
\textbf{是什么：}两断面之间有流体机械对流体做功时，能量方程要加一项：
水泵把机械能\textbf{输入}流体（加在方程左端），水轮机从流体取走机械能（加在下游端）。

\begin{gongshi}{带泵的伯努利方程与扬程}
\[ z_1+\frac{p_1}{\rho g}+\frac{\alpha_1 v_1^2}{2g}+H=z_2+\frac{p_2}{\rho g}+\frac{\alpha_2 v_2^2}{2g}+h_w \]
有能量输出（水轮机）时把左端的 $+H$ 换成右端的 $+H_{\text{水轮机}}$。\\
泵的\textbf{扬程} $H$（单位 m）：
\[ H=(z_2-z_1)+\frac{p_2-p_1}{\rho g}+\frac{\alpha_2 v_2^2-\alpha_1 v_1^2}{2g}+h_w \]
\textbf{含义：}\textbf{扬程 $H$＝泵使单位重力流体增加的能量}，即泵进出口断面的总水头之差。
液面敞开于大气、速度也为零时简化为 $H=(z_2-z_1)+h_w$——\textbf{扬程主要用于克服位差和水头损失}。
$H$ 的几何意义：若泵装在水面上、排出口竖直朝上且管路极短（损失可略），则液体能被扬起的高度就是 $H$。\\
\textbf{适用条件：}与总流伯努利方程的五条相同，只是两断面之间\textbf{存在能量输入/输出}；
水泵必须取\textbf{进口断面与出口断面}分别作为 1、2 断面，且 $H$ 对流体做正功。\\
\textbf{注意：}通过泵的流量称\textbf{泵的排量} $Q$；扬程不是"扬水高度"，还包括压能差、动能差和损失。
\end{gongshi}

\begin{gongshi}{泵的功率（有效功率、轴功率、电机功率）}
\[ N_{\text{有效}}=N_{\text{泵}}=\rho g Q H\ \ (\mathrm{W}),\qquad
   \eta_{\text{泵}}=\frac{N_{\text{有效}}}{N_{\text{轴}}},\qquad
   N_{\text{轴}}=\frac{\rho g Q H}{\eta_{\text{泵}}} \]
\[ N_{\text{电机}}=\frac{N_{\text{轴}}}{\eta_{\text{电机}}},\qquad
   \eta_{\text{总}}=\eta_{\text{泵}}\cdot\eta_{\text{电机}} \]
\textbf{含义：}$\rho gQH$＝单位时间内泵\textbf{传给流体}的能量，叫\textbf{有效功率（输出功率）}；
泵由轴输入功率，叫\textbf{轴功率}（也称额定功率）；轴功率再除以电机效率得到电机需要的功率。
$\eta_{\text{泵}}$＝泵效（有效功率/轴功率），$\eta<1$ 因为泵内部有水力、容积、机械三项损失。\\
\textbf{适用条件：}效率按题目给出的值代入；单位必须统一成国际单位（$Q$ 用 $\mathrm{m^3/s}$，
$H$ 用 m，$\rho$ 用 $\mathrm{kg/m^3}$，得 $N$ 为 W；题给 $\mathrm{kW}$ 要换算）。\\
\textbf{易混：}\textbf{扬程 $H$ 与功率是两个量}——同一台泵流量增大时扬程通常下降（泵的性能曲线），
不能由 $H$ 单独判断功率变化。
\end{gongshi}

\textbf{常考题型：}填空（"泵的扬程指（泵使单位重力流体增加的能量）"、"有效功率的表达式是（$\rho gQH$）"）；
计算（求扬程 $H$ → 求轴功率 $N_{\text{轴}}$ → 选电机；带泵供水/集水井/吸水管问题）；
问答（泵的工况点与节流调节：关小阀门使管路特性曲线变陡，流量减小、扬程略增，从而改变工况点）。
\textbf{重要度 \xstarfull{4}}\quad\yiju{E5 题库 \texttt{10.pdf} 计算题第 4 题（水泵工况点与节流调节）；E2 第四章"需掌握计算套路"；E6 教材 4.2 节含"泵对流体的作用"}\quad\nandu{中}

\begin{banfa}
\textbf{①带泵列方程的两个动作}：先把泵"切开"，取\textbf{泵进口断面为 1、出口断面为 2}；
再把 $H$ 写在\textbf{上游端（左端）加号}——因为泵给流体能量，下游总水头必高于上游。
水轮机则相反：$H_{\text{水轮机}}$ 写在下游端，因为流体把能量交给了机械。\\
\textbf{②"两断面之间"是判定 $h_w$ 的关键词}：$h_w$ 只算\textbf{取定两断面之间}那段管路的损失，
包括该段内所有沿程损失与局部损失；不要把管系里其他段的损失算进来。
题目没给损失数据（或说"损失忽略不计"）时 $h_w=0$。\\
\textbf{③功率计算的三步链条}：$(\rho gQH)\to$ 有效功率 $\to\div\eta_{\text{泵}}$ 得轴功率 $\to\div\eta_{\text{电机}}$ 得电机功率。
题目问"电机功率"却只除了一次效率是典型错误。\\
\textbf{④扬程的另一种求法}：若题给泵前后压力表的读数 $p_1$、$p_2$ 且两表同高，
则 $H=\dfrac{p_2-p_1}{\rho g}+\dfrac{v_2^2-v_1^2}{2g}$，比逐项列方程更快。\\
\textbf{⑤集水井/吸水管题的固定顺序}：先对"\textbf{井水面～泵进口}"列方程求出泵进口压强（必为负压），
再对"\textbf{泵进口～出水池}"列带泵方程求 $H$；吸水管题里最常考的量是\textbf{泵允许吸上真空高度}与最大安装高度。
\end{banfa}
\end{kaodian}

\begin{kaodian}{5}{伯努利方程的典型应用：文丘里管、皮托管、虹吸管、管嘴与孔口}
\textbf{是什么：}伯努利方程＋连续性方程的组合应用。核心套路是\textbf{把"求速度"的问题转化成"测压差"的问题}。

\textbf{（一）文丘里流量计（必考）}
由收敛段、喉道、扩散段组成，量流量用。取上游粗管断面 1 与喉道断面 2（都是渐变流断面），
忽略两断面间损失，与连续性方程 $v_1A_1=v_2A_2$ 联立：
\begin{gongshi}{文丘里管的流量公式}
\[ v_1=\frac{1}{\sqrt{\left(\dfrac{A_1}{A_2}\right)^2-1}}\sqrt{2g\!\left(\frac{p_1-p_2}{\rho g}+z_1-z_2\right)} \]
\[ Q=\mu A_1 v_1,\qquad \mu=\frac{1}{\sqrt{\left(\dfrac{A_1}{A_2}\right)^2-1}} \]
\textbf{含义：}流量与\textbf{两断面测压管水头差}的平方根成正比——这就是节流式流量计（孔板、喷嘴、文丘里管）
共同的原理："\textbf{流量与差压的平方根成正比}"。喉道处断面收缩、流速增大、压强降低，前后形成压差。\\
\textbf{适用条件：}不可压缩流体定常流动；两断面取在\textbf{渐变流处}（喉道是渐变流，可取值）；
两断面间损失略去或并入流量系数；实测时用流量系数 $\mu$（或 $\mu<1$）修正收缩与黏性影响。\\
\textbf{压差用 U 形水银压差计读}：$z_1-z_2+\dfrac{p_1-p_2}{\rho g}=\left(\dfrac{\rho_{Hg}}{\rho}-1\right)h$（$h$ 为水银柱差）。
\end{gongshi}

\textbf{（二）测速管（皮托管）}
一端开口正对来流，在管口形成\textbf{驻点}（流速为零），此处压强称驻点压强（总压）；
侧面另开小孔感受静压 $p$。对同一流线上的驻点与远前方未扰动点列伯努利方程，
两点等高、远前方流速即来流速度 $v$，驻点 $u=0$：
\begin{gongshi}{皮托管测速公式}
\[ v=\sqrt{\frac{2(p_0-p)}{\rho}}=\sqrt{2g\!\left(\frac{p_0}{\rho g}-\frac{p}{\rho g}\right)}=\sqrt{2gh} \]
\textbf{含义：}$p_0-p=\dfrac{1}{2}\rho v^2$ 即\textbf{动压}（流速水头对应的压强），
只要测出总压与静压之差就能算出流速。教材原图例中：B 点在上游未扰动处（$v_B=v$、压强 $p_B=\rho gH$），
A 点为驻点（$v_A=0$、压强 $p_A=\rho g(H+h)$），故 $v=\sqrt{2gh}$。\\
\textbf{适用条件：}不可压缩流体、定常流；探头轴线与来流方向平行（否则要乘方向修正）；
测速管测的是\textbf{某一点}的流速（点流），不是断面平均流速——要求 $v$ 还得靠流速分布或 $\alpha$。\\
\textbf{易混：}\textbf{皮托管测点流速、文丘里管测流量}，两者都靠"总压与静压之差"或"两断面压差"。
\end{gongshi}

\textbf{（三）虹吸管}：管路自高位水池翻越高处再下泄。取上游水面为 1、下游水面为 2（都取自由液面，
则 $p_1=p_2=0$ 表压、$v_1=v_2=0$），得\textbf{总水头差＝全部水头损失}（含进口、转弯、出口）。
虹吸管内\textbf{最高点压强一定小于大气压}（这是 2022 年单选与题库填空反复考的点）；
为防汽化，最高点负压不能超过\textbf{允许真空度}（水约 $7\sim8\ \mathrm{m}$ 水柱），
由此可反算\textbf{最大允许安装高度或虹吸高度}。

\textbf{（四）孔口与管嘴出流}（与第 5 章衔接）：取水池液面 1、出口断面 2，$\alpha=1$、$h_w$ 记为局部损失：
\begin{gongshi}{孔口/管嘴出流的流速与流量}
\[ v=\varphi\sqrt{2gH_0},\qquad Q=\mu A\sqrt{2gH_0},\qquad \mu=\varepsilon\varphi \]
\textbf{含义：}$H_0=H+\alpha v_0^2/(2g)$ 为作用水头（含行近流速水头）；$\varphi$ 为流速系数、$\varepsilon$ 为收缩系数、$\mu$ 为流量系数。
\textbf{适用条件：}定常出流、不可压缩；孔口自由出流时取收缩断面，管嘴出流取\textbf{管嘴出口断面}（此时 $\varepsilon=1$）。\\
\textbf{常考常数：}圆柱形直角外管嘴正常工作条件 $H_0\le 9\ \mathrm{m}$、$l=(3\sim4)d$；管嘴 $\mu\approx0.82$。
\end{gongshi}

\textbf{（五）常用节流装置对比}：
\begin{center}
\begin{tabularx}{\textwidth}{@{}lXX@{}}
\toprule
装置 & 结构特点 & 适用与注意 \\
\midrule
孔板 & 结构最简单，突然收缩→突然扩大，局部损失大 & 测气体与液体流量；流量系数需标定，永久压损大 \\
喷嘴 & 收缩段光滑，损失较孔板小 & 测蒸汽、气体；$\varepsilon$ 稳定 \\
文丘里管 & 由收敛段、喉道、扩散段组成，损失最小 & 精度高、造价高；喉道处压强低，注意防汽化 \\
\bottomrule
\end{tabularx}
\end{center}

\textbf{常考题型：}计算（文丘里管流量公式推导或求流量；管嘴体积流量；吸附/喷射泵求真空度）；
选择（2022 年单选考虹吸管最高处压强与大气压的关系；管嘴正常工作条件）；
填空（变径管求 $v_2$；U 形水银压差计读数换算）。
\textbf{重要度 \xstarfull{5}}\quad\yiju{E4 2015 年试题计算题含文丘里流量计、突然缩小水头损失；E4 2026 年回忆版计算题第 2 题（文丘里管流量证明）、第 5 题（管嘴体积流量，书上例题）；E4 2022 年单选第 7 题（虹吸管最高处压强）、第 8 题（圆柱形直角外管嘴正常工作条件）}\quad\nandu{中}

\begin{banfa}
\textbf{①文丘里"证明题"的三步模板}（2026 年真题就是这个题型）：
第 1 步 写出两个断面的伯努利方程（$z+\dfrac{p}{\rho g}+\dfrac{v^2}{2g}$ 相等）；
第 2 步 用连续性方程把 $v_2$ 换成 $v_1$（$A_1v_1=A_2v_2$），整理出 $v_1$ 的表达式；
第 3 步 乘断面面积得 $Q$，并说明 $Q\propto\sqrt{\Delta h}$。
证明题\textbf{不要求数字}，但每一步都要写明"对断面 1、2 列伯努利方程""由连续性方程 $A_1v_1=A_2v_2$"，
分步给分。\\
\textbf{②文丘里/孔板的差压用压差计读}：先用测压管水头差 $h=\left(\dfrac{p_1}{\rho g}+z_1\right)-\left(\dfrac{p_2}{\rho g}+z_2\right)$，
再按 $\left(\dfrac{\rho_0}{\rho}-1\right)h'$ 换算成水银柱差 $h'$，顺序颠倒数值必错。\\
\textbf{③"计算点"三选一速判}：\emph{自由液面}（$p=0$、$v\approx0$，未知数最少，首选）；
\emph{断面中心}（便于用 $z$ 直接读数）；\emph{驻点}（$v=0$，皮托管专用）。\\
\textbf{④管嘴题的出口断面取法}：外管嘴取\textbf{管嘴出口断面}（$\varepsilon=1$），
孔口则取\textbf{孔口后的收缩断面}。取错断面，流量系数就会用错。\\
\textbf{⑤虹吸管最大高度题套路}：先由上下游液面列方程求出\emph{总流量}（或由流量反求流速），
再对"上游液面～虹吸最高点"列方程求该点压强，令其等于允许真空度即得最大高度。
\end{banfa}

\begin{yicuo}
\textbf{①两断面压强基准混用}：一侧用绝对压强、另一侧用相对压强；$p_1=p_2$ 时看似"相等"，
实际相差一个大气压，结果必然错。\textbf{一律用相对压强（表压）}最省事。\\
\textbf{②断面取在急变流处}：弯管中、突扩/突缩的过渡段、闸门后都不能作为计算断面，
因为断面上不满足静压分布。\\
\textbf{③把 $h_w$ 加在上游侧}，或把"两断面之间"的损失算成了全管系的损失。\\
\textbf{④有能量输入/输出时不加 $H$}：题干出现"水泵""抽水站""水轮机"，必须换成带泵形式。\\
\textbf{⑤单位不统一}：$p$ 用 $\mathrm{kPa}$、$Q$ 用 $\mathrm{L/s}$、管径用 $\mathrm{mm}$ 时先全部换算成
$\mathrm{Pa}$、$\mathrm{m^3/s}$、$\mathrm{m}$ 再代入。\\
\textbf{⑥无旋与有旋不分}：题给"有旋流场"时，理想流体伯努利方程不能跨流线取两点（见 4.1）。
\end{yicuo}
\end{kaodian}

\section{动量方程}

\begin{kaodian}{5}{定常不可压缩总流的动量方程与分量式}
\textbf{是什么：}动量方程来自牛顿第二定律（动量定理）：\textbf{控制体内流体动量的时间变化率＝作用在控制体上的合外力}。
对定常、一元、只有一个进口和一个出口的控制体，流入与流出的动量差就等于合外力。
它解决的是伯努利方程\textbf{不能}解决的问题：\textbf{流体与固体边界的相互作用力}（弯管推力、射流冲击力、闸门受力）。

\begin{gongshi}{动量方程（矢量式与分量式，式 4.12）}
\[ \sum\bF=\rho Q\left(\beta_2\bv_2-\beta_1\bv_1\right) \]
\[
\begin{aligned}
\sum F_x&=\rho Q\left(\beta_2v_{2x}-\beta_1v_{1x}\right)\\
\sum F_y&=\rho Q\left(\beta_2v_{2y}-\beta_1v_{1y}\right)\\
\sum F_z&=\rho Q\left(\beta_2v_{2z}-\beta_1v_{1z}\right)
\end{aligned}
\]
多个进出口时：
\[ \sum\bF=\sum_{\text{出}}\rho Q_i\bv_i-\sum_{\text{进}}\rho Q_i\bv_i \]
\textbf{含义：}方程\textbf{左端是外力，右端是动量的变化}。右端＝单位时间内\textbf{流出}控制体的动量
减去\textbf{流入}控制体的动量，即\textbf{动量的净输出率}。
$\sum\bF$ 是作用在\textbf{控制体内流体}上的\textbf{全部外力}，只包含三类：
\begin{xiaowen}
  \item \textbf{断面压力}：$p_1A_1$、$p_2A_2$，方向\textbf{由外向内}压向控制体（指向控制体内部）；
  \item \textbf{重力}：$G=\rho g V$（控制体内流体的重量），竖直向下；水平弯管问题中重力与流动平面垂直，分量常为 0；
  \item \textbf{固体壁面对流体的作用力} $R$：\emph{待求量}。题目问的"流体对弯管的作用力"是它的\textbf{反作用力} $-R$。
\end{xiaowen}
\textbf{适用条件（必背）：}① 定常流动；② 不可压缩流体；③ 所取控制体的两断面都是\textbf{渐变流断面}
（这样才能用静压分布计算断面压力 $pA$）；④ 只计外力（质量力＋表面力＋约束反力），不计内力
（控制体内部流体之间的相互作用力成对抵消）；⑤ 一元流动近似（用断面平均流速乘 $\beta$）。
\textbf{动量修正系数 $\beta$：}按点流速算的实际动量与按平均流速算的假想动量之比，
$\beta=\dfrac{\displaystyle\int_A u^2\,\dd A}{v^2 A}\ \ge 1$，反映断面速度分布不均匀性。
层流（抛物线分布）$\beta=\dfrac{4}{3}$；湍流 $\beta\approx 1$。题目不说明则取 $\beta=1$。\\
\textbf{与 $\alpha$ 的记法区别：}$\alpha$ 出现在\emph{能量}方程里、是\textbf{三次方}平均（层流 $2$）；
$\beta$ 出现在\emph{动量}方程里、是\textbf{二次方}平均（层流 $4/3$）。
\end{gongshi}

\begin{tishi}
\textbf{用动量方程还是伯努利方程？}\ \textbf{求力}（弯管对流体/流体对弯管的作用力、射流冲击力、闸门受力）
$\to$ \textbf{动量方程}；\textbf{求压强、流速、流量、水头损失、功率} $\to$ \textbf{伯努利方程}。
两者常需\textbf{联立}：伯努利方程给出两断面压强差，动量方程再把压强差变成力。
这是 4.3 计算题的标准结构。
\end{tishi}

\textbf{常考题型：}计算（2015 年试题计算题：水平弯管用动量方程求 $F_x$、$F_y$；2026 年回忆版计算题第 3 题＝书上动量方程例题 4.12）；
选择/判断（"动量方程是标量式还是矢量式"、"动量方程中压强用相对压强还是绝对压强"）；
填空（动量修正系数的层流取值 $4/3$）。
\textbf{重要度 \xstarfull{5}}\quad\yiju{E4 2015 年试题计算题考水平弯管动量方程求 $F_x$、$F_y$；E4 2026 年回忆版计算题第 3 题为书上动量方程例题 4.12；E1 slide16 计算考动量方程；E1 slide33 第四章＝考试计算题重点}\quad\nandu{难}

\begin{banfa}
\textbf{①"动量的变化在右侧"绝不放反}：$\sum\bF=\rho Q(\beta_2\bv_2-\beta_1\bv_1)$。
记忆句：\textbf{"外力＝出减进"}。若写成 $\rho Q(\bv_1-\bv_2)$，得到的是壁面\emph{对流体}的反号，
即是"流体对壁面"的力——\textbf{题目问的对象不同，符号就不同}，一定要在最后一步明确回答的是谁对谁的力。\\
\textbf{②坐标系的取法}：一般令 $x$ 轴沿\textbf{出口}（或上游主流）方向，
把速度投影写全（$v_{1x}=v_1$、$v_{1y}=0$），能减少符号错误。斜板/曲面射流题必须画方向余弦
（$v_{2x}=v_2\cos\theta$、$v_{2y}=v_2\sin\theta$）。\\
\textbf{③断面压力用相对压强}：$p_1A_1$、$p_2A_2$ 中的 $p$ 必须与方程里其他压强同一基准；
用表压时，断面压力也写成表压 $\times A$。\\
\textbf{④外力"漏项自检"}：三个问题——\emph{断面压力两项写了吗？重力写了吗？未知的壁面反力设了吗？}
水平弯管中重力与水平面垂直，$x$、$y$ 分量都为 0，但竖直弯管或变径竖直段就必须计入。\\
\textbf{⑤未知数配平}：动量方程本身只能求力；若 $v_1$、$v_2$、$p_1$、$p_2$ 有多个未知，
先用\textbf{连续性方程}定 $v_2$，再用\textbf{伯努利方程}定 $p_2$，最后代入动量方程求力。
\textbf{方程数（连续性＋伯努利＋动量）≥ 未知量个数}才可解。\\
\textbf{⑥动量矩方程}（见 4.4）只在涉及\emph{旋转机械或力矩}时用，一般弯管题不必列力矩式。
\end{banfa}
\end{kaodian}

\begin{kaodian}{5}{动量方程的典型题：水平弯管、变径弯管、射流冲击、喷嘴}
\textbf{是什么：}动量方程的四类经典应用，都遵循同一套"取控制体→选断面→建坐标→列分量式→联立连续性/伯努利"流程。

\textbf{（一）水平弯管（2015 年真题题型）}
等径水平弯管，断面 1、2 都在渐变流处，$v_1=v_2=v$，两断面压力 $p_1$、$p_2$ 由伯努利方程或题给数据确定：
\[
\sum F_x=p_1A_1-p_2A_2\cos\theta+R_x=\rho Q\left(v_2\cos\theta-v_1\right)
\]
\[
\sum F_y=-p_2A_2\sin\theta+R_y=\rho Q\left(v_2\sin\theta-0\right)
\]
解出 $R_x$、$R_y$，则\textbf{弯管（约束）对水流的力}为 $\bR=(R_x,R_y)$，
\textbf{水流对弯管的作用力}为 $-\bR$，大小 $|\bR|=\sqrt{R_x^2+R_y^2}$，方向 $\tan\phi=R_y/R_x$。
水平弯管中重力沿 $z$ 方向，$x$、$y$ 分量均为 0，故 \textbf{重力不进入水平分量式}。

\textbf{（二）变径弯管}：$A_1\ne A_2$，先由连续性方程得 $v_2=v_1A_1/A_2$，
再由伯努利方程求 $p_2$，然后照（一）列分量式。此时动量的两个分量一般都不能忽略。
竖直管段或竖弯管必须把重力 $\rho gV$ 计入相应分量。

\textbf{（三）射流冲击固定平板}：射流从喷嘴以 $v$ 冲击平板，控制体取喷嘴出口到射流沿平板流出处。
$z$ 向动量变化为 $0-\rho Qv$（取射流方向为正），故
\[ F=\rho Qv=\rho A v^2 \]
斜板（与射流成 $\theta$）时分力 $F_n=\rho Qv\sin\theta$；射流分成两支，用$x$、$y$ 分量式与流量分配联立。
曲面（如半圆形叶片）射流反转 $180^\circ$：$F=\rho Q(v-(-v))=2\rho Qv$（\textbf{倍力}现象，水轮机动力的来源）。
碰撞后速度保留系数为 $k$ 时把 $v$ 换成 $kv$。

\textbf{（四）喷嘴与弯管联立}：水流经收缩喷嘴加速喷出，喷嘴受的拉力由动量方程求得
$F=\rho Q(v_1-v_2)$（配合断面压力 $p_1A_1$ 使用）；求喷嘴处压强则由伯努利方程得到。
这就是"\textbf{用动量方程求力、用伯努利方程求压强}"的两方程联用结构。

\textbf{常考题型：}计算（2015 年试题水平弯管求 $F_x$、$F_y$；2026 年回忆版第 3 题＝书上例题 4.12——均为弯管类）；
选择（射流冲击平板的力与流速的几次方成正比——$\rho Av^2$，\textbf{二次方}）。
\textbf{重要度 \xstarfull{5}}\quad\yiju{E4 2015 年试题计算题（水平弯管动量方程求 $F_x$、$F_y$）；E4 2026 年回忆版计算题第 3 题＝书上动量方程例题 4.12；E2 第四章"需掌握计算套路"}\quad\nandu{难}

\begin{zhenti}{2026 年 818 计算题第 3 题（回忆版）}{计算题}
\textbf{书上动量方程例题 4.12}（教材 4.3 节动量方程例题）。\\
\textbf{题型判定：}弯管（变径弯管）类动量方程综合题——由连续性方程求出口流速、由伯努利方程求出口压强、
再由动量方程求水流对弯管（或弯管对水流）的作用力 $F_x$、$F_y$ 及其合力；若断面在竖直方向，
还要把重力计入。参考答案要点见下方"弯管动量方程五步"。
\end{zhenti}
\noindent 本题为 2015 年"水平弯管动量方程求 $F_x$、$F_y$"的同类改编，\textbf{大概率原题或改数字}。
教材例题的数值随版本印次可能不同，考试按题目所给数据现场列式即可，\textbf{套路完全一致}。

\begin{banfa}
\textbf{弯管动量方程五步（背下来就能做任何弯管题）：}
\begin{ti}
  \item \textbf{取控制体}：把弯管内两断面之间的水体整体取出，画出它受到的力（两个断面压力、重力、壁面反力 $R$）。
  \item \textbf{选断面}：断面 1、2 都必须取在弯曲段\textbf{上游和下游的直管渐变流段}，
        不能取在弯头内部（否则不能用静压分布算 $pA$）。
  \item \textbf{建坐标}：以弯管所在平面为 $xOy$，$x$ 轴沿入口（或出口）方向，
        把 $v_1$、$v_2$ 与 $p_1A_1$、$p_2A_2$ 全部分解到坐标轴上，避免"忘记分解压力"。
  \item \textbf{列分量式}：$x$、$y$ 各一个方程，注意右端永远是"出减进"，
        且壁面反力 $R_x$、$R_y$ 先按\textbf{正方向假设}，解出负数说明方向与假设相反（要写出来）。
  \item \textbf{联立补齐}：$v_2$ 由连续性方程得、$p_2$ 由伯努利方程得，
        最后求合力大小与方向，并\textbf{明确回答题目问的是"流体对弯管的力"（＝$-R$）还是"弯管对流体的力"（＝$R$）}。
\end{ti}
\textbf{补充技巧：}\ 等径水平弯管的合力可用几何法速算——两断面压力与动量变化都是矢量，
画矢量三角形 $p_1A_1+p_2A_2+(-\rho Q\bv_1)+\rho Q\bv_2+R=0$（控制体受力封闭），
比逐个分量更直观；但\emph{书写求解过程时仍要写分量式}，否则步骤分拿不到。
\end{banfa}

\begin{yicuo}
\textbf{①动量方程是矢量式，必须投影}：直接写 $\sum F=\rho Q(v_2-v_1)$ 而上代入矢量，
在弯管、射流题中一定错。\textbf{分量式才是可计算的方程}。\\
\textbf{②外力要算全}：断面压力（两个）、重力、固体壁面反力，一个都不能漏；
最常漏的是\textbf{出口断面的压力 $p_2A_2$}（"出口排到大气"时 $p_2=0$ 表压，但\emph{有压管出口}不是大气）。\\
\textbf{③压强用相对压强时，断面压力 $pA$ 也必须用相对压强}：
若方程里混用绝对压强，$p_1A_1$ 的量级会多出 $p_aA$ 的巨大项，结果荒谬。\\
\textbf{④方向符号搞反}：$\rho Q(\bv_2-\bv_1)$ 是"动量净输出"，等于\emph{外力}的合力；
求"流体对固体的力"要再取反号。\\
\textbf{⑤忽略 $\beta$ 或在层流题里取 1}：层流 \textbf{$\beta=4/3$}，题目若给出层流条件应按 $4/3$ 计，
不给则取 1；与此对应能量方程中层流 \textbf{$\alpha=2$}——两个系数不要互换。
\end{yicuo}
\end{kaodian}

\section{动量矩方程}

\begin{kaodian}{2}{动量矩方程与离心式流体机械的欧拉方程}
\textbf{是什么：}把动量定理用于\textbf{矩}——对某个固定点（或轴）取矩，得到动量矩方程。
它描述的是流体力矩与流体动量矩变化之间的关系，是叶片式流体机械（水泵、水轮机、风机）的基本方程来源。

\begin{gongshi}{定常流动的动量矩方程}
\[ \sum\bM=\rho Q\left(\beta_2\,\br_2\times\bv_2-\beta_1\,\br_1\times\bv_1\right) \]
绕 $z$ 轴的分量式：
\[ \sum M_z=\rho Q\left(\beta_2 r_{2}\,v_{2u}-\beta_1 r_{1}\,v_{1u}\right) \]
\textbf{含义：}作用在控制体上的\textbf{合外力矩}＝单位时间内流出与流入控制体的\textbf{动量矩之差}。
$r$ 为计算点到给定轴的距离（力臂），$v_u$ 为速度在\textbf{圆周（切向）方向的分量}。
含义与动量方程平行：把"力"换成"力矩"、把"动量 $\rho Qv$"换成"动量矩 $\rho Qrv_u$"。\\
\textbf{适用条件：}与动量方程相同——定常、不可压缩、断面取渐变流、只计外力（矩）；
此外必须\textbf{先指定取矩的轴}，所有力臂都对该轴量取。\\
\textbf{常考形式：}判断/填空（"动量矩方程是矢量式还是标量式"、"动量矩方程用于求\textbf{力矩}"）。
\end{gongshi}

\begin{gongshi}{离心式流体机械的欧拉方程（简述）}
叶轮进出口断面取半径 $r_1$、$r_2$ 处的圆周面，$u$ 为牵连（圆周）速度、$v_u$ 为绝对速度的切向分量：
\[ H_T=\frac{\omega}{g}\left(r_2v_{2u}-r_1v_{1u}\right)=\frac{u_2v_{2u}-u_1v_{1u}}{g} \]
\[ M=\rho Q\left(r_2v_{2u}-r_1v_{1u}\right) \]
\textbf{含义：}泵（水轮机）的理论扬程与叶轮进出口的\textbf{速度矩} $rv_u$ 之差成正比，
功率 $N=M\omega=\rho gQH_T$。物理上就是\textbf{叶轮把机械能通过动量矩的变化传给流体}。\\
\textbf{适用条件：}理想（无黏）流体、叶片数无穷多、流体沿叶片相对运动为定常；
实际还要用水力效率、容积效率修正。教材把该部分放在流体机械章节，初试只作概念了解。
\end{gongshi}

\textbf{常考题型：}\textbf{真题中本章 4.4 未见独立成题}，属了解内容；
只需能说出"动量矩方程用来求流体对机械（叶轮）的力矩"、"离心机械欧拉方程联系扬程与速度矩 $rv_u$"。
\textbf{重要度 \xstarfull{2}}\quad\yiju{E1 slide33 第四章＝考试计算题重点（章节级依据）；E2 第四章"需掌握计算套路"（章节级依据）；E4 历年真题中 4.4 节未见独立成题}\quad\nandu{易}

\begin{banfa}
\textbf{①与动量方程对照记}：动量方程 $\to$ 求\emph{力}（$\rho Qv$）；动量矩方程 $\to$ 求\emph{力矩}
（$\rho Qrv_u$）。两者的右端都是"出减进"，左端都是"外力/外力矩"，
只是把线量换成角量。\\
\textbf{②$\bv_2$ 要分解}：叶片机械题里绝对速度 $\bv$ 必须分成\textbf{径向分量}（过流量，
决定 $Q$ 与连续性）和\textbf{切向分量 $v_u$}（决定动量矩与扬程），轴向分量对轴无矩。\\
\textbf{③考试取舍}：本章 4.4 只背 $\sum M_z=\rho Q(r_2v_{2u}-r_1v_{1u})$ 与
$H_T=\dfrac{u_2v_{2u}-u_1v_{1u}}{g}$ 两个形状，看懂"速度矩"一词即可，
\textbf{不要在 4.4 上花时间}——把时间投到 4.2 与 4.3 的计算套路上更划算。
\end{banfa}
\end{kaodian}

\section{本章小结}

\begin{center}
\begin{tabularx}{\textwidth}{@{}lccX@{}}
\toprule
考点 & 重要度 & 难度 & 一句话抓住 \\
\midrule
理想流体特征与欧拉方程 & \xstarfull{3} & 易 & $\bm{f}-\frac{1}{\rho}\grad p=\frac{\dd\bu}{\dd t}$；只要求理想流体，适用范围最广 \\
兰姆--葛罗米柯形式 & \xstarfull{3} & 易 & 加速度 $=\grad(u^2/2)-\bu\times\bomega$；认出"旋涡项"即可 \\
理想流体伯努利积分 & \xstarfull{3} & 易 & 四条适用条件；有旋时只沿流线、无旋才全场 \\
实际流体总流伯努利方程 & \xstarfull{5} & 中 & 加 $\alpha$、加 $h_w$；两断面必须同一基准、取渐变流断面 \\
带泵/水轮机与泵功率 & \xstarfull{4} & 中 & 左端 $+H$；$N_{\text{轴}}=\rho gQH/\eta$，除两次效率 \\
典型应用（文丘里/皮托管/虹吸管/管嘴） & \xstarfull{5} & 中 & 伯努利＋连续性；流量与差压的平方根成正比 \\
动量方程与分量式 & \xstarfull{5} & 难 & $\sum\bF=\rho Q(\beta_2\bv_2-\beta_1\bv_1)$，外力＝出减进 \\
弯管/射流/喷嘴典型题 & \xstarfull{5} & 难 & 取控制体→两渐变流断面→建坐标→列分量式→联立连续性＋伯努利 \\
动量矩方程与欧拉方程 & \xstarfull{2} & 易 & 求力矩；$H_T=(u_2v_{2u}-u_1v_{1u})/g$，了解一下即可 \\
\bottomrule
\end{tabularx}
\end{center}

\noindent\textbf{本章计算题固定套路（合二为一）}：
伯努利方程管\textbf{能量（压强、流速、损失、功率）}，动量方程管\textbf{力（约束反力、冲击力）}，
连续性方程管\textbf{流量（速度与断面面积的换算）}。
一道大题通常就是"连续性定速度 $\to$ 伯努利定压强 $\to$ 动量定力"这三步，
把这三步写成固定动作，本章的 15--20 分就稳了。
